\chapter{Auditoría de fuentes y criterio editorial}

\section{Fuentes proporcionadas}
\begin{longtable}{p{.28\textwidth}p{.64\textwidth}}
\caption{Fuentes proporcionadas y función que cumplen en esta edición.}
\label{tab:auditoria-fuentes}\\
\toprule
\textbf{Fuente} & \textbf{Uso en esta edición}\\
\midrule\endhead
Guía docente 2026--2027 & Fija el temario oficial, competencias, evaluación y bibliografía recomendada.\autocite{ule-guia-2026}\\
Transparencias Temas 1--2 (partes 1 y 2) & Hipótesis básicas, concepto de tensión, cargas, vínculos, criterio de signos, esfuerzos internos y ecuaciones de equilibrio interno.\autocite{lanza-temas12a,lanza-temas12b}\\
Introducción de estructuras aeronáuticas & Contextualización de fuselaje, ala, largueros, larguerillos, costillas y semimonocasco.\autocite{ule-estructuras-intro}\\
Corrección ordinaria enero 2025 (Wuolah) & Referencia de estilo de examen: reacciones/leyes de esfuerzos y progresión hacia problemas de sección, cortante y torsión.\autocite{wuolah-examen-2025}\\
Página pública de Wuolah para la asignatura & Confirma la existencia de material de Teoría de Estructuras de 3.º de Aeroespacial ULe, incluidos exámenes/correcciones 2025 y 2026.\autocite{wuolah-ule-page}\\
\bottomrule
\end{longtable}

\section{Fuentes abiertas de contraste}
MIT OpenCourseWare ofrece el texto completo de \textit{Engineering Mechanics for Structures} y materiales de mecánica estructural; TU Delft ofrece un curso específico de \textit{Aerospace Mechanics of Materials}, con teoría, ejemplos y ejercicios sobre tensión, deformación, axil, torsión, flexión y cortante.\autocite{bucciarelli2002,tudelft-mom,mit-structural-mechanics,mit-analysis-design}

Para resistencia de materiales en español se utiliza también el texto de Cervera y Blanco publicado por CIMNE, con ISBN verificable y acceso desde la página académica del autor.\autocite{cervera2015}

\section{Conflictos y correcciones}
\begin{warningbox}
Los apuntes del curso emplean una conversión aproximada $1\,\mathrm{MPa}\approx10\,\mathrm{kp/cm^2}$. La equivalencia más precisa es $1\,\mathrm{MPa}\approx10.197\,\mathrm{kp/cm^2}$. En cálculos modernos se recomienda mantener SI.
\end{warningbox}

Las convenciones de signo pueden variar entre bibliografías. En este volumen prevalece el convenio explícito del curso para $N$, $V$ y $M$; cuando se usa una fuente externa se traduce antes de combinar fórmulas.

\section{Lagunas detectadas}
El material detallado proporcionado se concentra en el Bloque I. La guía docente confirma temas posteriores de flexión/axil, torsión de pared delgada, cortante, energía y elasticidad tridimensional, pero esos bloques no están documentados con el mismo nivel de detalle en los PDFs proporcionados. Por eso esta edición se identifica de forma explícita como \textbf{Volumen I}; ampliar el libro sin esa advertencia mezclaría material oficial detallado con reconstrucción externa de forma poco transparente.
